% ============================================================
% Related Work section -- meant to be \input{} or \include{}'d from a
% main manuscript file, not compiled standalone. Same bibliography/style
% requirements as introduction.tex (see that file's header comment).
%
% CITATION STATUS: every \cite{...} below resolves to a real, verified
% entry in references.bib. Entries added specifically for this section
% (verified against the publisher page / PubMed / Crossref):
%   - Sun:2024 (npj Digital Medicine review of PD digital biomarkers
%     broadly -- imaging, biofluid, wearables -- used in 2.1 to situate
%     voice/speech among the wider objective-biomarker landscape)
%   - Eliwa:2025 (PeerJ Comput. Sci. -- PSO jointly tuning acoustic
%     feature selection and classifier hyperparameters for PD prediction)
%   - Appati:2026 (Discover Applied Sciences -- chaotic Grey Wolf-Dragonfly
%     hybrid for voice-driven PD prediction in high-dimensional data)
%   - Ali:2023 (Diagnostics -- filter FS + genetic algorithm + ensemble
%     learning on voice-only PD datasets)
%   These three broaden 2.3's opening citation list beyond the five
%   source papers this study reproduces as baselines, evidencing that
%   evolutionary/swarm-based FS for PD voice data is an active, current
%   line of work rather than represented only by this paper's own
%   baseline choices.
% ============================================================

\section{Related Work}
\label{sec:2}

\subsection{Current Diagnostic Practice and the Case for Objective Biomarkers}
\label{sec:2.1}

Clinical diagnosis of Parkinson's disease (PD) remains anchored to the MDS
Clinical Diagnostic Criteria \cite{Postuma:2015}: a judgment, made by a
trained examiner, of whether bradykinesia is present together with rest
tremor and/or rigidity, refined by supportive features and exclusion
criteria and typically corroborated by the patient's response to levodopa
therapy. Because this judgment rests on the visual and manual assessment
of motor signs that can be subtle, intermittent, or altogether absent
early in the disease course, diagnostic accuracy is known to vary
considerably with examiner experience and disease duration, and
misdiagnosis remains common enough at early stages to be a recurring
concern in the movement-disorders literature \cite{Rizzo:2016}. This has
motivated a broader search for objective biomarkers that do not depend on
a clinician's subjective interpretation of motor signs. Dopaminergic
imaging (DAT-SPECT, PET), cerebrospinal-fluid and blood-based
$\alpha$-synuclein assays, genetic testing, and, more recently, wearable
and other digital-sensing modalities have all been investigated as
candidate objective markers \cite{Sun:2024}; each can support diagnosis or
staging, but each also carries some combination of high cost, invasiveness,
specialized equipment, or limited accessibility outside tertiary care
centers \cite{Sun:2024}. Voice and speech recordings, by contrast, can be
acquired with nothing more than a standard microphone, at negligible
marginal cost, and without any physical contact with the patient --- a
combination of properties that makes them an unusually attractive
candidate among the emerging digital-biomarker modalities, and the
specific focus of this study.

\subsection{Acoustic and Speech Biomarkers of Parkinsonian Dysphonia}
\label{sec:2.2}

Hypokinetic dysarthria, characterized by reduced vocal loudness,
monopitch, imprecise articulation, and breathy or hoarse voice quality, is
estimated to affect the large majority of individuals with PD at some
point in the disease course \cite{Ho:1999,Cao:2025}. Its physiological
basis lies in rigidity and bradykinesia affecting the laryngeal,
respiratory, and articulatory musculature: reduced vocal-fold adduction
force and control lead to incomplete or irregular glottal closure, reduced
respiratory support lowers subglottal pressure and vocal intensity, and
reduced range of motion in the articulators blurs consonant and vowel
targets \cite{Yang:2020}. These physiological changes are reflected in
several well-established acoustic correlates: increased cycle-to-cycle
perturbation in fundamental frequency (jitter) and amplitude (shimmer),
reduced harmonics-to-noise ratio (HNR) reflecting incomplete glottal
closure, reduced vowel space area reflecting imprecise articulation, and
altered nonlinear dynamical measures of vocal-fold vibration -- including
recurrence period density entropy (RPDE), detrended fluctuation analysis
(DFA), and pitch period entropy (PPE) -- that are more sensitive than
classical perturbation measures to the aperiodic, weakly chaotic component
of parkinsonian phonation \cite{Yang:2020}. Because these physiological
changes can begin before overt motor signs are clinically apparent,
acoustic and broader speech biomarkers have increasingly been proposed as
tools for prodromal detection and longitudinal monitoring, not only
post-diagnosis screening \cite{Rusz:2024,Cao:2025}. This physiological
grounding is revisited in Section~\ref{sec:discussion} when interpreting
which acoustic features the proposed pipeline selects most consistently.

\subsection{Metaheuristic Feature Selection for PD Detection}
\label{sec:2.3}

Wrapper-based feature selection couples a search procedure with a
downstream classifier's cross-validated performance as the objective to
optimize. Population-based metaheuristics -- derivative-free and
requiring no assumption of a smooth or convex fitness landscape -- have
become a standard choice of search procedure for this task across
biomedical applications generally \cite{BarreraGarcia:2023}. Within
voice-based PD detection specifically, a fast-growing line of work applies
swarm- and evolution-inspired algorithms to select acoustic feature
subsets ahead of a downstream classifier, frequently reporting accuracies
above 90\% on public benchmarks such as the Oxford voice dataset:
particle-swarm-, butterfly-, and whale-inspired variants
\cite{Hashemi:2026}, particle swarm jointly tuning the feature subset and
classifier hyperparameters \cite{Eliwa:2025}, a modified grey wolf
optimizer \cite{Santhosh:2025}, a modified hunger games search
\cite{Hashim:2023}, a quantum mayfly optimizer \cite{Mansour:2024}, a
genetic algorithm with filter-based pre-selection \cite{Ali:2023}, a
beluga whale optimizer \cite{Ganesan:2025}, and hybrids chaining two
algorithms in sequence -- a whale-to-grey-wolf cascade
\cite{AlNajjar:2024} and a chaotic grey wolf-dragonfly cascade
\cite{Appati:2026}. As discussed in Section~\ref{sec:introduction}, the
classifier consuming the selected features is fixed and pre-determined in
the large majority of this work, rarely optimized jointly with the
feature subset within a single search process -- the gap this study's
proposed encoding is designed to close, by searching a substantially
higher-dimensional, joint feature-and-classifier space without
sacrificing search quality relative to feature-selection-only
formulations. The optimizer underlying this study's proposed method, the
Competitive Swarm Optimizer (CSO) \cite{Cheng:2015}, departs from the
canonical PSO-style update rule in which every particle is pulled toward a
single global-best position. Instead, CSO randomly pairs particles each
generation into one-on-one competitions: the fitter particle (the winner)
of each pair passes to the next generation unchanged, while the weaker
particle (the loser) updates its position and velocity toward the winner,
with an additional pull toward the swarm's mean position. Because no
particle is ever pulled directly toward a single global best, CSO was
originally designed to preserve population diversity and resist the
premature convergence that standard PSO exhibits on large-scale problems
-- a property this study leverages directly, since the joint
feature-and-classifier encoding is considerably larger than a
feature-selection-only encoding of the same dataset, making resistance to
premature convergence practically relevant rather than only a
large-scale-benchmark curiosity. Section~\ref{sec:3} describes the
three strategies this study adds on top of plain CSO -- elite-guided
update, stagnation-triggered opposition-based learning, and a
diversity-preserving elite archive -- together forming the proposed
Elite-guided Opposition-based Archive Competitive Swarm Optimizer
(EOACSO).
